\documentclass[book.tex]{subfiles}

\begin{document}
\chapter{Kernel Methods}

\label{chap:kernel_methods}
\noindent\rule{11cm}{0.4pt} \\
\textbf{Prerequisites:}  \autoref{chap:ml_basics}, \autoref{chap:hilbert},  \\
\textbf{Difficulty Level:} ** \\
\noindent\rule{11cm}{0.4pt}
\vspace{5mm}

Kernel methods offer a structured method of adding nonlinear features to linear regression. The basic intuition behind a kernel method is to compare a given query point $x$ against known datapoints $x_i$ to determine prediction $y$. For example, suppose that we have the training dataset

\begin{align}
\mathcal{D} &= \{(x_i, y_i) \}
\end{align}

Then to make a prediction at query point $x$, we predict
\begin{align}
    y &= \sum_i K(x, x_i) y_i
\end{align}
here $K$ is some distance function capturing the notion of distance between $x$ and $x_i$. Points closer to $x_i$ will make predictions closer to $y_i$. Many possible families of kernel functions $K$ exist. Note that in the case $K$ only depends on $x$ and not $x_i$, kernel methods are exactly the same as linear regression.

\section{Radial Basis Function Kernels}

The radial basis function (RBF) kernels are a broadly used family of nonlinear kernels

\begin{align}
K(x, x') &= \exp \left ( - \frac{\|x - x'\|^2}{2 \sigma^2} \right )
\end{align}
Here $\sigma$ is a hyperparameter for the kernel. RBF kernels are closely related to Gaussians, and arise in many physical systems.

\begin{lstlisting}
def rbf($X: \mathbb{R}^{M \times N}$, $\sigma$) $\to$ $\mathbb{R}^{M \times M}$:
  K: $\mathbb{R}^{M \times M}$
  for i j
    K[i,j] = exp(-$\|X[i]-X[j]\|^2$/($2\sigma^2$))
  return K
\end{lstlisting}

\section{Efficiency Considerations}

Kernel methods require the computation of $K$ for every datapoint in the training dataset. For large datasets, the expense can become prohibitive. For this reason, a vast literature of approximate kernel methods have emerged. These methods typically attempt to factor the matrix $K(x_i, x_j)$ into smaller matrices and use these matrix factorizations to allow for faster query answers.

\section{Exercises}
\begin{enumerate}
\item Prove that linear regression is a special case of kernel regression for a suitable kernel $K$.
\end{enumerate}
\end{document}